\section{Analýza konkurenčních řešení}

% Write - hned sem na začátek napsat jaké jsou cíle analýzy - že chci odhalit jakým způsobem se vyučuje softwarové inženýrství a jaké jsou příležitosti pro zlepšení této výuky, jaké silné a slabé stránky aplikace mají, abychom mohli vytvořit aplikaci, která vyplňuje tyto nedostatky a pak zjistit "inspiraci" (to asi nazvat jinak) ohledně dalších funkcí a nástrojů, které existující aplikace poskytují

% Write - Napsat sem někam podle čeho jsem aplikace k analýze vybíral. Měl bych to mít dobře odůvodněné.

% Write
V rámci analýzy konkurenčních řešení jsem se rozhodl zvolit nejpoužívanější aplikace na výuku programování a softwarového inženýrství, které se blíží vizi tohoto projektu.

% Write - zde sjednotit vizi s úvodem (tohle udělat až budu mít Done úvod)
Zaměřím se tedy zejména na aplikace, které jsou určeny pro samouky a kladou důraz na kvalitní uživatelskou zkušenost a mikrolearning. % Write - zde ověřit, jestli existuje mikrolearning



tedy aplikace, které budou přímou. Mezi tyto konkurenty patří aplikace Mimo, CodeAcademy, Scrimba a ???.  

% Write
V analýze se zaměřím nejprve na vzhled uživatelského rozhraní a uživatelskou zkušenost. Poznatky budou klíčové při návrhu nových funkcí a nového vzhledu uživatelského rozhraní. Následně se zaměřím na způsoby učení, motivace uživatelů a způsoby monetizace aplikace. Tyto informace budou sloužit k formulaci požadavků tohoto projektu. Dále se v analýze zaměřím na silné a slabé stránky jednotlivých aplikací a na další poznatky, které mohou být použity při návrhu této aplikace.
